\documentclass[a4paper,13pt,titlepage,oneside]{article}
\usepackage[greek,spanish]{babel}
\usepackage{setspace}
\usepackage[utf8]{inputenc}
\usepackage{times}
\usepackage{graphicx}
\usepackage{apacite}
\usepackage{csquotes}
\usepackage[T1]{fontenc}
\usepackage{anyfontsize}
\usepackage{fancyhdr}
\pagestyle{fancy}

%\fancyhf{}
%\lhead[\leftmark]{Marco Debiasi}
%\rhead[Constantinopla en el siglo XI]{\rightmark}
%\lfoot[\thepage]{}
%\rfoot[]{\thepage} 

\bibliographystyle{apacite}

\setcounter{secnumdepth}{0}
%\setcounter{tocdepth}{1}

\newcounter{ns}
\addtocounter{ns}{1}

%\title{Constantinopla en el siglo XI}
%\author{Marco Patricio Debiasi}
%\date{\today}
%\pagestyle{myheadings}
%\markright{Marco P. Debiasi}

\begin{document}
	\begin{titlepage}
	\centering{\fontsize{40pt}{48pt}
	\selectfont Constantinopla en el siglo XI \par}
	\vspace{1.5cm}

	\centering{\fontsize{22pt}{20pt}\selectfont\includegraphics[width=0.7\linewidth]{ico.png}}

	\fontsize{22pt}{20pt}\selectfont Marco Patricio Debiasi\par
	\vspace{1cm}
	\fontsize{22pt}{10pt}\selectfont Profesorado de Historia\par
	\end{titlepage}
\spacing{1.5}


\newpage

\section{Marco de Referencia}

%El contexto científico en el que se va a encontrar el trabajo, 
%se trata de recopilar información de los diversos estudios que 
%hablen del tema elegido. Ya que el tema elegido (cómo se vivía en la 
%Constantinopla del siglo XI) no encontre nadie que trate de 
%forma directa el ``cómo se vivía'' como título o 
%interrogante principal, los antecedentes del estudio son los que 
%investigaron brindando la información necesaria para la reconstrucción 
%de cómo pudo ser la vida en la Constantinopla del siglo XI, y los 
%novelistas, que estudiando y leyendo para la construcción histórica 
%del escenario de sus novelas, estudiaron el tema de forma directa, 
%como es el objetivo del trabajo.


%Primer revisión, cambio de parrafo
Me propongo analizar la vista cultural, social, militar y educativa de
Constantinopla durante el siglo XI. Elegí ese siglo por
la lectura del libro de la Alexiada en 2019, citada en el
presente trabajo.
En este trabajo me propongo recopilar información de diversos
estudios en los que se aborda el tema elegido.
Esas investigaciones me brindan la información necesaria para la
reconstrucción de cómo pudo ser la vida en la Constantinopla
del siglo XI. A esto se agrega, un trabajo que pueda servirme
y pueda servirle a novelistas para la construcción del escenario
de sus novelas.

%Nuevo parrafo, tras revisión 1
El tema central del trabajo es la reconstrucción de la vida
en dicha ciudad y tiempo, tal como se enuncia en la primer
interrogante.

Interrogantes e hipótesis del trabajo:

¿Cómo se vivía en la Constantinopla del siglo XI?

%Parentesis a la interrogante
¿Cuál era el plan de estudios?
(De la Universidad de Constantinopla, y la preocupación
por la educación del emperador Alejo (último libro))

¿Cómo era la salud pública?

¿Cómo se organizaba el ejército?

¿Quiénes eran los personajes político-militares-culturales?

A pesar de que varias de las anteriores interrogantes ya las habrán 
tratado de forma directa, en este trabajo se recopilaría dicha 
información con el fin de reconstruir la vida de Constantinopla del 
siglo XI, como figura anteriormente.

%Agregado autor
Por ejemplo, en el libro XV de “La Alexiada”, 
de Ana Comneno traducido por Emilio Díaz
Ronaldo, se relata un gran hospital 
público, y un colegio, publico, para huérfanos, junto con una nueva 
didáctica de dicho colegio. Además, se relata una nueva formación militar.

%Agregado autor
Y, en el capítulo IV de “Filosofía Bizantina”, 
por Basilio Tatakis, se describe el plan de 
estudios que había diseñado el catedrático de la universidad de 
Constantinopla, Miguel Psellos.

El método que se va a usar en el trabajo, es el método bibliográfico. 
Ya que se va a confiar en los testimonios de los que estudiaron el 
tema, y en la fuente primaria de “La Alexiada”.



\newpage

\section{Introducción}
	En el presente trabajo se desarrollará varios aspectos de cómo pudo haber sido la vida en la Constantinopla del siglo XI; viendo su filosofía, cultura, educación, salud y política. Cuya organización está hecha para que cada capítulo sea independiente de los demás, para que se pueda leer como si fueran un conjunto de artículos.

\newpage

\section{Capítulo 1: Filosofía Bizantina}
\subsection{La Mística}
	\subsubsection{San Simeón, el nuevo teólogo}
	\input{filosofia/san_simeon}
	\subsubsection{Calixto Kataphigiotis}
	\input{filosofia/calixto}
\subsection{La Moral}
	\input{filosofia/moral}
	\subsubsection{Cecaumeno}
	\subsubsection{Nikoulitzas}
	\subsubsection{Teofilacto de Bulgaria}
	\subsubsection{Neoplatonismo}
	\input{filosofia/neoplatonismo}


\newpage

\section{Capítulo 2: Cultura}
\subsection{}
\newpage

\section{Capítulo 3: Educación}

%Para empezar con la educación de Constantinopla, se va a dividir en 
%dos partes principales, el rector de la universidad de Constantinopla, 
%y el emperador del imperio. De forma que se va a poder estudiar el 
%programa de estudio y que obra hizo el “estado”.

%Modificación de la primer corrección
Se va a dividir en dos partes, en la primera voy a estudiar
la figura del rector de la Universidad de Constantinopla
Miguel Psellos; y en la segunda la relación entre el emperador y la
educación.

Voy a empezar a desarrollar por la parte pedagógica, con la obra e 
ideas de Miguel Psellos, continuando por el emperador Comneno, 
ya que gobernaron durante el siglo XI. En el primero, se va a dividir 
en una presentación del personaje junto con su visión educativa y 
programa de materias. En el segundo, se mencionará la preocupación por 
la educación del emperador, del colegio de huérfanos.

\subsection{Miguel Psellos}
En la presente sección, desarrollaré cómo era la educación en la 
Constantinopla del siglo XI, permaneciendo fiel al objetivo del 
presente trabajo de averiguar y detallar cómo era la vida en este 
determinado espacio-tiempo.

\subsubsection{Presentación}
Para empezar, relataré brevemente la historia de este personaje. 
Miguel Psellos vivió desde 1018 hasta el 1096 después de cristo, 
admirador de la cultura helénica y de Platón, aunque más viendo era 
neoplatónico, siendo \textit{“[...]el importador de este peculiar neoplatonismo 
en el pensamiento bizantino [...]”.}
\footnote{\cite[p.~160]{filosofia}}

\subsubsection{Visión Pedagógica}
Primero veremos el método de enseñanza o pedagogía de Psellos; 
\textit{“[...]el método de la enseñanza que daba en la Universidad de 
Constantinopla; hallaremos así indicios característicos de la idea de 
Psellos tenía de la filosofía y que refleja cuando dice: 
“Antes de dedicarme al estudio de la filosofía, estuve enamorado de 
la retórica”. […] Con ciertas variantes expresa en otras ocasiones 
la misma idea, como por ejemplo:“Yo combino la filosofía con la 
retórica y procuro expresarme de acuerdo con ambas y con la ayuda de 
cada una de ellas.” En otros términos: que reclama armonía de 
formas y contenido. “El cuidado del estilo -insiste- no es, 
en modo alguno, obstáculo para el mérito.””}
\footnote{\cite[p.~160]{filosofia}}

Para ir viendo el estilo de pedagogía que quería dar en la 
Universidad de Constantinopla, también: \textit{“[...]La retórica ejerce 
sobre Psellos una especie de fascinación, por lo que dice: 
“La sustancia del retórico no es más que el resplandor del lenguaje, 
se que hable, o que escriba, o que enseñe la retórica, 
o que se ocupe su pensamiento con la filosofía…; esto es lo que 
hace que el retórico no sea ni ateniense ni espartano, sino ciudadano 
del universo; es como el relámpago y el trueno, cuyo fulgor y sonido 
se imponen por doquier.””.}
\footnote{\cite[p.~161]{filosofia}}

Entrando más en el tema pedagógico de Miguel Psellos, continua el autor:

\textit{“A pesar de su fascinación a la retórico, Psellos no se decide 
exclusivamente por ella al tratar del problema pedagógico que separaba 
a los profesores de retórica (o sofistas) de los filósofos. 
“Es preciso -escribe a cierto padre- que pongas a tu hijo en medio de 
las fuentes, la filosofía y la retórica, y que le hagas beber de ambas, 
de cada una oportunamente; si no lo procuras así, le convertirás en 
una inteligencia desprovista de lenguaje, si es que sólo bebe de la 
fuente filosófica; o resultará mera lengua sin inteligencia en el 
caso contrario.”}
\footnote{\cite[p.~161]{filosofia}}

Aunque, tampoco -dice- alcanzan para formar a un hombre, sino que 
necesita también de política \textit{“[...]si no, ese hombre, con toda su 
retórica sólo será un 
	\selectlanguage{greek}
	χύμβαλον άλαλάζον.
	\selectlanguage{spanish}
 Destacamos así un rasgo característico de Psellos, ya que ni la 
 filosofía ni la retórica son para él un fin en sí, sino meros 
 instrumentos, de los que se sirve para educar a la juventud bizantina, 
 a la que ha sabido inculcar su amor a Grecia.”}
 \footnote{\cite[p.~161]{filosofia}}

Por lo cual, puede verse el interés de él por ser un buen 
orador para enseñar.

\textit{“Según otro punto de vista suyo, el valor de la filosofía y de la 
retórica es muy relativo, si no nulo, ya que comparados con la 
Sagrada Escritura, los retóricos y los filósofos, 
los caldeos y los egipcios, todos son como el 
bronce comparado con el oro”.}
\footnote{\cite[p.~161]{filosofia}}

\subsubsection{Programa}
%Modificación revisión 1
En cuanto lo visto anteriormente, Psellos desarrollo su 
programa para la Universidad de Constantinopla; donde seguramente
estudiaban retórica, dialéctica, son tres disciplinas antes
de pasar a Quadrivium;
el trivium, aritmética, geometría, 
música, astronomía y el quadrivium.

\textit{“Fiel a estas ideas, Psellos estableció en la Universidad de 
Constantinopla el programa siguiente: Primero se instruía a los 
principiantes en retórica y en dialéctica; terminado el trivium, 
los estudiantes pasaban al curso superior, el de Psellos, en el que 
estudiaban aritmética, geometría, música, astronomía, las ciencias 
enumeradas en el libro IV de la República de Platón y el quadrivium. 
Pero aquí no pasaban a la dialéctica (como quería Platón, ni en el 
sentido que él preconizaba), sino que seguían con los estudios de 
filosofía, por considerar a ésta como el complemento de todas las 
ciencias. Desde luego, el estudio de la filosofía 
(que se comenzaba por la lógica de Aristóteles) no constituía un 
nuevo grado, sino que era más bien la continuación y culminación 
del quadrivium.”}
\footnote{\cite[pp.~161--162]{filosofia}}

Dentro de la universidad, la cátedra de filología la tenía un tal 
Nicetas, el cual usaba un tal método de ciencias para la 
gramática proveniente de Alejandría:

\textit{“Conviene señalar que el curso de gramática comprendía la 
filología entra; el estudio de los diversos dialectos griegos 
hacía resaltar más las leyes principales de esta lengua; la 
ley de analogía, que se aplicaba sobre todo para explicar la 
declinación, ayudaba a comprender la lengua en su regularidad 
racional. Las teorías rítmicas y musicales venían a completar el 
estudio del lenguaje. En resumen: que Nicetas usaba ampliamente la 
ciencia de los gramáticos de Alejandría:”}
\footnote{\cite[p.~162]{filosofia}}

Miguel Psellos, era un personaje religioso, que el autor lo marca 
lejos de tener una ciencia positiva, al mirar las ciencias y 
mitología griega desde las Sagradas Escrituras.


\subsubsection{Retórica}
Ya que hemos empezado a adentrándonos en la filología, gramática y 
retórica anteriormente, con los métodos de enseñanza de Psellos y 
Nicetas, continuemos con las tres principales materias del estudio, 
la retórica/oratoria, la filosofía y la metafísica.

A continuación, continuaré con oratoria:

\textit{“No solo Nicetas de ocupaba de la filología, pues también Psellos 
enseñaba gramática, si bien sus preferencias tendían al examen de 
las leyes que rigen el mundo de la oratoria, procurando que sus discípulos 
adquirieran cualidades estilísticas. Su examen teórico iban seguidos 
de ejercicios prácticos, más lecturas y explicaciones de páginas 
de oradores y de poetas, elegidas unas y otras por la perfección de 
su forma y el interés de su contenido. Se cuenta que, para tal examen, 
Psellos seguía a Hermógenes de Tarso; pero aunque no hay duda de 
que conocía los escritos de éste y se servía de ellos, su objetivo era 
(según declara formalmente) enseñar la retórica antigua, 
“la que el propio Platón prefería”, una retórica que no gustaba a 
sus discípulos (quienes preferían la de Hermógenes) y para la cual 
se guiaba siempre por Platón y Aristóteles.”}
\footnote{\cite[p.~163]{filosofia}}

\subsubsection{Quadrivium}
En cuanto a este curso, el autor relata los cursos 
que se enseñaban en el quadrivium.

\textit{“Para enseñar las ciencias del quadrivium, Psellos utilizó los 
manuales de Nicómaco de Gerasa, de Euclides, de Diofante y de 
Teón de Esmirna para las matemáticas, los de Tolomeo y Proclo para la 
astronomía y los de Aristógenes para la música. 
Al tratar de las ciencias físicas y naturales, insiste particularmente 
en las teorías de lo frío y de lo cálido, lo seco y lo húmedo, 
la explicación de los fenómenos meteorológicos y de los terremotos, 
así como la descripción de la Tierra. 
También se incluye en sus enseñanzas la adivinación, 
la astrología y las narraciones maravillosas, 
las que pretende explicar aclarando que los hechos ocurridos en esos 
dominios no son extraordinarios sino en apariencia, 
ya que en su fondo están regidos por las leyes científicas."}
\footnote{\cite[p.~163]{filosofia}}

\subsubsection{Filosofía}
Este era un curso, dividido en lógica y metafísica.

\textit{“El objetivo del curso de filosofía era dar, mediante la lógica y 
la metafísica aristotélicas, los fundamentos del pensamiento 
filosófico, a la vez que servía de punto de partida para el 
estudio de los problemas especulativos. Consideraba a la filosofía 
como una etapa preparatoria de la metafísica, cuyo material era tomado 
casi íntegramente de Plotino, Proclo y Platón, y no de Aristóteles. 
Las doctrinas filosóficas de la metafísica tenían a la teología, 
filosofía primera, como referencia máxima, y a la luz de ella debían 
los estudiantes interpretar los textos teológicos. 
Tras esto se puede comprender en todo su valor el sentido de la 
definición de Psellos daba a sus lecciones, de las que: 
“La lección es la perfección del alma, su redacción y su ascensión, 
o su vuelta al bien supremo."}
\footnote{\cite[pp.~164--165]{filosofia}}

\subsubsection{Metafísica}
En cuanto a la última materia que desarrollá Tatakis, 
es la Metafísica, en la cual pone;

\textit{“Además de la filosofía, constituían asignaturas de la etapa 
preparatoria de la metafísica la historia de la filosofía, 
la exégesis de las narraciones legendarias y de las tradiciones 
y sentencias populares… siempre que no se utilizan como argumentos 
para probar la verdad de los resultados adquiridos. 
La historia de la filosofía no se limitaba  a relatar el pensamiento 
griego, sino que abarcaba también las doctrinas provenientes del 
Oriente (de Caldea, Egipto y Palestina) e inclusive el pensamiento 
hermético. 
Por otra parte, el análisis de las tradiciones populares y de las 
narraciones legendarias de la Antigüedad conseguía extraer de ellas 
ideas sublimen que Psellos trasladaba al cristianismo. 
Este programa de estudios muestra el papel importantísimo, central, 
que se le había reservado a la tradición helénica; 
de ella se esperaba no sólo que modelara el justo y el juicio de la 
juventud, sino que adornará además su alma con valores morales; 
y notemos bien que Psellos procura mostrar esta tradición helenística 
tanto en sus aciertos como en sus imperfecciones, 
considerándola siempre como precursora del cristianismo.”}
\footnote{\cite[p.~165]{filosofia}}

\subsection{Ana y Alejo Comneno}


\subsubsection{El orfanato}

% Introducción, orfanato-escuela
\textit{``
	Dedicó el día siguiente entero al cuidado de
	cautivos y recién llegados. Repartió entre todos
	aquellos de sus allegados, que sabía que llevaban una
	vida honesta, y entre los higúmenos [abades] de los sagrados
	monasterios a todos los niños que habían quedados
	privados de padres y estaban sumidos en la amarga desgracia
	de la orfandad, y les recomendó que no los criasen como
	esclavos, sino como seres libres, considerándolos merecederos
	de una completa formación e intruyéndolos en las
	Sagradas Escrituras. También entregó algunos al orfanato que
	él había fundado y que estaba pensado más como escuela para
	quienes quisieran aprender, a fin de que sus directores
	les enseñeran el ciclo completo de estudios [trivium y quadrivium].
	''}
\footnote{\cite[pp.~605--606]{alexiadaXV}}

Para concluir el capítulo sobre educación, presento la narración sobre
el orfanato-escuela fundado por el emperador según la autora.
Dónde describe brevemente el lugar, menciona quienes van, el método
didáctico que se empezó a usar ahí.

% Descripción, educación en el orfanato
\textit{``
	[...]. A la derecha del gran templo [de los apóstoles] hay una escuela
	primaria para los niños huérfanos procedentes de toda variedad de razas,
	en donde un maestro imparte la clase y los niños se colocan en torno
	a él, unos atemorizados por las preguntas sobre grámatica, otros
	escribiendo la denominada esquedografía. Allí es posible ver a un
	latino que se está instruyendo, a un escita que aprende griego, a un
	romano manejando textos griegos y a un griego iletrado que aprende hablar
	correctamente griego. Esos eran los afanes de Alejo sobre la formación
	intelectual. En cuanto a la técnica de la esquedografía, diremos que es
	un invento de los más recientes y originario de nuestra generación. [...]
	ahora el estudio de estos maestros, de los poetas, de los historiadores
	y de sus cualidades no ocupa siquiera un lugar secundario. El único interés
	es el juego, los demás trabajos están prohibidos.
	''}
\footnote{\cite[p.~609]{alexiadaXV}}

Por lo que se puede ver que no se necesitaba de mucho dinero para acceder
a una educación completa (trivium y quadrivium), además de poder leer
sobre el método didáctico que se estaba poniendo a prueba en dicho momento.

























\newpage

\section{Capítulo 4: Salud}

Antes de empezar con el tema de salud, me gustaría aclarar que
%lo siguiente es sacado de
para desarrollar este aspecto me he basado en
la fuente primaria del 
presente trabajo, “La Alexiada”, del libro XV, el último de la obra. 
%En el que se narra
% Segunda Revisión
En este se relatan
las obras que había mandado hacer el emperador, 
según lo narra su hija, para atender las necesidades de
los huérfanos y demás gente afectada 
por las continuas guerras que estaba teniendo el imperio.

Este es el más corto de los capítulos, ya que en la fuente
que estaré utilizando, mencionada abajo, no son muchas
páginas ni mucho contenido que hable al respecto.

En el libro XV, los párrafos del 4 al 8 están dedicados a 
describir un hospital que el emperador había mandado a construir, 
siendo el único fragmento destinado a la salud dentro del libro 
(sin contar a la sanidad de campaña (militar) que se describe 
en el mismo libro XV, sección VII, párrafos 1 y 2).

En dicha cita, la autora narra la última campaña militar
del emperador, dónde describe el trato con `civiles' y
enfermos que tranportaban. Una vez llega a la ciudad,
se describe la obra de salud y educación.
Dónde, la parte de educación está en el capítulo ononimo,
mientras que prosigo con la descripción sobre salud.

\subsection{Salud de campaña}
En la Alexiada, se relata cómo marchaba el ejército, en la
última campaña militar del emperador, y hacía atender
a la gente que llevaba:

\textit{``Cuando llegó a Filomelio, tras liberar por doquier
	a quienes estaban bajo el yugo bárbaro [...], introdujo
	en medio de la formación a cautivos, mujeres, niños y
	todo el botín y se puso en el camino de vuelta apaciblemente
	con un movimiento lento y semejante al de las hormigas.
	''}

\textit{``
	Como había muchas mujeres embarazadas y otras muchas sufrían
	enfermedades, cuando a una mujer estaba a punto de dar a luz,
	tocaba la trompeta a una señal del soberano, todos se detenían
	enseguida y la formación entera se quedaba quieta en el mismo
	lugar. Nada más enterarse de que el parto había concluido,
	mandaba a dar otro toque, que no era de los habituales,
	comunicando la puesta en marcha y animando a todos a caminar.
	Si alguien se estaba muriendo, de nuevo sucedía lo mismo y el
	soberano se presentaba en el sitio donde yacía el moribundo,
	llamaba a los sacerdotes para que cantaran los himnos postreros
	y para que le dieran los sacramentos al agonizante y, una vez se
	habían celebrado las honras fúnebres a los difuntos de acuerdo
	con las normas sagradas, hasta que el muerto no estuviera enterrado,
	no permitía que se moviera lo más mínimo la formación. A la hora de comer,
	hacía llamar a todas las mujeres y hombres que estuvieran agotados
	por las enfermedades y la vejez, les ofrecía lo mejor de su comida
	y ordenaba a sus comensales que hicieran lo mismo.
	''}\footnote{\cite[pp.~604--605]{alexiadaXV}}

Quise poner este breve relato ya que cuenta el trato con los enfermos
y `civiles' que transportaba con su ejército.

\subsection{Salud en Constantinopla}

En cuanto a la salud en Constantinopla, cuando el emperador llega
a la ciudad, manda hacer un orfanato para los hijos de los soldados
caídos, y se reponsabiliza por la educación (que está en el capítulo
sobre educación en la parte sobre el emperador).

Y, junto con el orfanato para los hijos de los caídos, levantaría,
cómo relata en la fuente, un gran hospital para enfermos, incluyendo
salas de `internación'.

\textit{``
	En el sector que existe junto a la acrópolis, donde se abre
	el acceso al mar, había encontrado un templo de enorme tamaño
	bajo la advocación del gran apóstol Pablo, y construyo allí,
	dentro de la ciudad imperial, otra ciudad. [...] En su interior
	hay erigidas circularmente un conjunto abigarrado de viviendas,
	moradas para los más pobres y, lo que demuestra mayor caridad,
	hospicios para personas mutiladas.
	''}

\textit{``
	Este recinto circular es doble y gemelo. Los unos, hombre y
	mujeres mutilados, habitan en la parte superior. Otros se arratran
	en la planta baja.
	''}\footnote{\cite[p.~606]{alexiadaXV}}

La autora relata su experiencia viendo el lugar:

\textit{``
	Yo misma he llegado a ver a una mujer vieja asistida por una
	joven, a un hombre ciego guiado por manos de uno que sí ve,
	a personas sin pies que poseían pies, no lo suyos propios
	sino los de otros; a personas sin manos auxiliadas por las de
	otras personas, a criaturas recién nacidas amamantadas por otras
	madres, a paralíticos servidos por otros hombres robustos.
	Era doble la muchedumbre que recibía alimentos, pues unos se
	contaban entre los servidos y otros entre los servidores.
	''}\footnote{\cite[p.~607]{alexiadaXV}}

Finalizando con una descripción del lugar:

\textit{``
	Al templo del gran apóstol Pablo lo dotó con un importante y
	numeroso clero y con abundancia de lámparas. Si se visita este
	templo, se puede ver cómo cantan dos coros, uno a cada lado,
	alternativamente. Pues, como hizo Salomón, dispuso la existencia
	en el templo de los apóstoles de cantantes masculinos y femeninos.
	''}\footnote{\cite[p.~608]{alexiadaXV}}



































\newpage

\section{Capítulo 5: Política}
\subsection{Relaciones con Italia}
Para empezar a relatar sobre las relaciones políticas que tenía Bizancio,
me gustaría hacer un recuento de las fuentes que dispongo para analizar
su política y diplomacia.

Empiezo por la fuente primaria “La Alexiada”, 
la fuente que inició este trabajo de investigación, 
dónde se relatan diversas campañas y alianzas 
que pueden darnos un gran panorama, 
por último, un libro del historiador “Roger Crowley”, 
para sustentar parte de lo que dice la fuente primaria. 
Pero, primordialmente, 
se va a poner a “La Alexiada” como, prácticamente, 
única fuente del presente capítulo del trabajo de investigación.

Dentro de la fuente primaria,  Ana Comneno, “La Alexiada”, 
Ático de los libros, Barcelona, 2016, en las páginas 249 a la 253, 
se puede leer la alianza con Venecia (que más adelante provocaría 
la guerra entre Venecia y Génova que ahogarían a Bizancio), luego, 
en las páginas 373 a la 291 y 419 a la 460, se puede leer la 
primera cruzada, y por ende, la opinión de los romanos 
orientales de la misma. Cómo para resaltar, en las páginas 
215 a la 218, se hace referencia a la querella de las investiduras, 
pero la autora desconoce la causa real (por lo menos eso parece) 
y cree que la causa de que Alemania (Sacro Imperio) invada al 
papa es otra… Por último, las invasiones normandas, 
turcas y de las tribus del nómadas. 

Finalizando este pequeño resumen, pantallazo, empecemos a 
relatar sobre cada uno de estos asuntos para conocer el 
panorama político y diplomático del Imperio Romano de Oriente.

Empecemos por la alianza Bizantino-veneciana, y la primer 
invasión normanda. Para entender el contexto, en los libros 
IV, V y VII de la Alexiada, se relata la primer invasión normanda 
(en el libro XII la segunda), el la cual Roberto Guiscardo, 
el duque de Apulia y futuro rey de Sicilia busca nuevas tierras de 
conquista en Grecia, dando inició a la guerra. Aquí se puede ver el 
primer aspecto diplomático, hay una guerra entre el imperio bizantino 
y sicilia. En esta guerra, el emperador del momento, Alejo, 
como lo narra la fuente, llama a Venecia (que prácticamente 
era como una colonia de Bizancio al comienzo de su historia), 
la cual derrota en el mar a los normandos.

“Roberto se hizo cargo de toda la flota [normanda] y, 
[…] desde donde zarparía en dirección al Ilírico.[…]

Roberto se percató de la ofensiva que pretendía la escuadra en contra 
de él y, anticipándose a la batalla de acuerdo con su carácter, soltó 
amarras y con toda su flota arribó al puerto Casope. Los venecianos, 
a su vez, llegaron al puerto de Pasaron y aguardaron allí un cierto 
tiempo. Cuando se enteraron de la llegada de Roberto, marcharon 
rápidamente también ellos al puerto de Casope. Tras un violento 
combate y un enfrentamiento al abordaje, Roberto fue derrotado. 
No por ello se rindió después de esta derrota, habida cuenta de su 
temperamento belicoso y dispuesto para el combate, sino que de nuevo 
se preparaba para luchar en otra batalla y enfrentarse en un 
combate más trascendente. Al conocer esto, los comandantes de ambas 
flotas, animados por la victoria, lo atacaron tres días después y 
lograron una brillante victoria sobre él. Luego, regresaron de 
nuevo al puerto de Pasaron.[…]

Transcurrido un tiempo, los venecianos aparejaron dromonos, 
tirremes y algunas otras naves pequeñas y veloces, y se 
encaminaron con mayores fuerzas contra Roberto.”
\footnote{\cite[pp.~250--252]{alexiadaVI}}

“Él [el emperador de Bizancio] les correspondió con abundantes presentes y
honores. Honró al dux de Venecia con la dignidad de protosebasto junto con 
sus rentas. Honró también al patriarca con la dignidad de hipértimo en 
unión de sus correspondientes rentas. Igualmente, ordenó que anualmente 
fuera distribuida entre todas las iglesias de Venecia una importante 
cantidad de oro procedente del tesoro imperial. Hizo tributarios a todos los 
naturales de Melfi que poseyeran negocios en Constantinopla de la 
iglesia del apóstol evangelista San Marcos y cedió la explotación de 
los negocios que se extendían desde el antiguo muelle de los hebreos 
hasta el lugar llamado Bigla, incluidos los muelles existentes dentro 
de estos límites. Les regaló asimismo muchos inmuebles en la ciudad 
imperial, en Dirraquio y en donde se les antojase pedirlos. Y, lo que 
es más importante, les concedió la exención de aranceles en el comercio 
dentro de las fronteras del Imperio de los romanos, para que comerciasen 
libremente a voluntad, sin tener que aportar ni un óbolo en virtud de tasas 
comerciales o de cualquier clase de impuesto exigido para los fondos 
públicos, así como la dispensa de subordinarse a ninguna autoridad romana.”
\footnote{\cite[p.~252]{alexiadaVI}}

El historiador Roger Crowley expone brevemente la relación 
que tenían estos dos:

“La relación entre Bizancio y Venecia fue intensa, compleja y longeva; 
también fue complicada, debido a sus visiones contradictorias del mundo, 
y estuvo sujeta a brutales cambios de humor. Sin embargo, Venecia siempre 
miró hacia Constantinopla. Era la gran ciudad del mundo, la puerta de 
entrada a oriente. A través de sus almacenes en el cuerno de oro fluía 
la riqueza del ancho mundo: pieles, cera, esclavos y caviar de rusia; 
especias de la India y China, marfil, seda, piedras preciosas y oro. 
A partir de estos materiales los artesanos bizantinos creaban objetos 
extraordinarios, tanto sagrados como profanos -relicarios, mosaicos, 
cálices con esmeraldas incrustadas o vestidos de seda tornasolada- 
que conformaban el gusto veneciano. La asombrosa basílica de San Marcos, 
re consagrada en 1094, fue diseñada por arquitectos griegos siguiendo la 
pauta de la iglesia matriz de los santos apóstoles en Constantinopla; 
sus artesanos relataron la historia de San Marcos, piedra a piedra, 
imitando el estilo de los mosaicos de Santa Sofía; sus Orfebres y 
esmaltadores crearon la Pala D’Oro, el retablo dorado, una expresión 
milagrosa de la devoción y el arte bizantinos. 
El aroma de las especias de los muelles de Venecia había sido 
transportado mil millas, desde los almacenes del cuerno de oro. 
Constantinopla era el zoco de Venecia, donde sus mercaderes se reunían y
 ganaban (o perdían) fortunas. Cómo leales súbditos del emperador, 
 el derecho a comerciar en sus tierras fue siempre su posesión 
 más preciada. El imperio, a su vez, lo utilizó como elemento de 
 negociación para mantener a raya a sus arrogantes vasallos. 
 En 991, Orseolo consiguió valiosos derechos de comercio a cambio 
 de apoyo veneciano en el Adriático, veinticinco años después, 
 estos derechos le fueron retirados con enfado tras una disputa.”
 \footnote{\cite[pp.~39--40]{venecia}}


\subsection{La querella de las envestiduras}

\subsection{Las tribus del norte del Danuvio}

\subsection{La primera cruzada}
\newpage

\section{Conclusión}
\subsection{Conclusión}
\newpage

\section{Bibliografía}
\bibliography{biblio}
\newpage

\setstretch{1.5}
\tableofcontents

\end{document}

